\documentclass[11pt]{article}

% Change "review" to "final" to generate the final (sometimes called camera-ready) version.
% Change to "preprint" to generate a non-anonymous version with page numbers.
\usepackage[review]{acl}

% Standard package includes
\usepackage{times}
\usepackage{latexsym}

% For proper rendering and hyphenation of words containing Latin characters (including in bib files)
\usepackage[T1]{fontenc}
% For Vietnamese characters
% \usepackage[T5]{fontenc}
% See https://www.latex-project.org/help/documentation/encguide.pdf for other character sets

% This assumes your files are encoded as UTF8
\usepackage[utf8]{inputenc}

% This is not strictly necessary, and may be commented out,
% but it will improve the layout of the manuscript,
% and will typically save some space.
\usepackage{microtype}

% This is also not strictly necessary, and may be commented out.
% However, it will improve the aesthetics of text in
% the typewriter font.
\usepackage{inconsolata}

%Including images in your LaTeX document requires adding
%additional package(s)
\usepackage{graphicx}
\usepackage{amsmath}

% If the title and author information does not fit in the area allocated, uncomment the following
%
%\setlength\titlebox{<dim>}
%
% and set <dim> to something 5cm or larger.

\title{Instructions for *ACL Proceedings}

% Author information can be set in various styles:
% For several authors from the same institution:
% \author{Author 1 \and ... \and Author n \\
%         Address line \\ ... \\ Address line}
% if the names do not fit well on one line use
%         Author 1 \\ {\bf Author 2} \\ ... \\ {\bf Author n} \\
% For authors from different institutions:
% \author{Author 1 \\ Address line \\  ... \\ Address line
%         \And  ... \And
%         Author n \\ Address line \\ ... \\ Address line}
% To start a separate ``row'' of authors use \AND, as in
% \author{Author 1 \\ Address line \\  ... \\ Address line
%         \AND
%         Author 2 \\ Address line \\ ... \\ Address line \And
%         Author 3 \\ Address line \\ ... \\ Address line}

\author{First Author \\
  Affiliation / Address line 1 \\
  Affiliation / Address line 2 \\
  Affiliation / Address line 3 \\
  \texttt{email@domain} \\\And
  Second Author \\
  Affiliation / Address line 1 \\
  Affiliation / Address line 2 \\
  Affiliation / Address line 3 \\
  \texttt{email@domain} \\}

%\author{
%  \textbf{First Author\textsuperscript{1}},
%  \textbf{Second Author\textsuperscript{1,2}},
%  \textbf{Third T. Author\textsuperscript{1}},
%  \textbf{Fourth Author\textsuperscript{1}},
%\\
%  \textbf{Fifth Author\textsuperscript{1,2}},
%  \textbf{Sixth Author\textsuperscript{1}},
%  \textbf{Seventh Author\textsuperscript{1}},
%  \textbf{Eighth Author \textsuperscript{1,2,3,4}},
%\\
%  \textbf{Ninth Author\textsuperscript{1}},
%  \textbf{Tenth Author\textsuperscript{1}},
%  \textbf{Eleventh E. Author\textsuperscript{1,2,3,4,5}},
%  \textbf{Twelfth Author\textsuperscript{1}},
%\\
%  \textbf{Thirteenth Author\textsuperscript{3}},
%  \textbf{Fourteenth F. Author\textsuperscript{2,4}},
%  \textbf{Fifteenth Author\textsuperscript{1}},
%  \textbf{Sixteenth Author\textsuperscript{1}},
%\\
%  \textbf{Seventeenth S. Author\textsuperscript{4,5}},
%  \textbf{Eighteenth Author\textsuperscript{3,4}},
%  \textbf{Nineteenth N. Author\textsuperscript{2,5}},
%  \textbf{Twentieth Author\textsuperscript{1}}
%\\
%\\
%  \textsuperscript{1}Affiliation 1,
%  \textsuperscript{2}Affiliation 2,
%  \textsuperscript{3}Affiliation 3,
%  \textsuperscript{4}Affiliation 4,
%  \textsuperscript{5}Affiliation 5
%\\
%  \small{
%    \textbf{Correspondence:} \href{mailto:email@domain}{email@domain}
%  }
%}

\begin{document}
\maketitle
\begin{abstract}
This document is a supplement to the general instructions for *ACL authors. It contains instructions for using the \LaTeX{} style files for ACL conferences.
The document itself conforms to its own specifications, and is therefore an example of what your manuscript should look like.
These instructions should be used both for papers submitted for review and for final versions of accepted papers.
\end{abstract}

\section{Introduction}
\label{sec:introduction}

% TODO: Introduction -- not yet written.

\section{Methods: OntoDisco}
\label{sec:methods}

% \subsection{Overview}
% \label{sec:methods-overview}

OntoDisco turns a corpus of open-vocabulary into a canonicalized, typed knowledge graph: a deduplicated relation vocabulary, a deduplicated and hierarchically organized entity-type vocabulary, a deduplicated entity vocabulary with class assignments, and a set of relation-level domain/range constraints.
At first, a schema-free LLM extractor  (\S\ref{sec:extraction}) produces a collection of raw triplets: 
\begin{equation}
\mathcal{G} = (s, t_s, r, o, t_o, Q),
\end{equation}
where $s$, -- subject entity, $t_s$ -- subject type, $r$ -- relation predicate, $o$ -- object entity, $t_o$ -- object type, $Q$ -- an optional list of qualifiers extracted from texts. Then $\mathcal{G}$ is passed further for canonization and typing. Every downstream stage resolves fields into canonical identifiers, keeping each stage's data dependencies explicit and independently inspectable.

The pipeline runs ordered stages, described in
\S\ref{sec:relation-dedup}--\S\ref{sec:constraints}: \emph{(1) relation
canonicalization}, \emph{(2) entity-type canonicalization}, which also
induces a taxonomy hierarchy over the resulting type vocabulary,
\emph{(3) entity deduplication and class assignment}, and \emph{(4)
domain/range constraint induction}. The ordering is dictated by data
dependencies. Entity-type canonicalization is conditioned on the canonical relation vocabulary as a disambiguating signal. Entity deduplication runs after the type hierarchy because it uses the hierarchy as a symbolic gate for candidate synonyms. Constraint induction derives relation constraints directly from the extracted triplets and the outputs of the first two stages.

The following design choices recur across stages. (1) Deduplication of synonimous surface forms is performed by clustering with Contriever sentence embeddings \citep{izacard2022contriever} and HDBSCAN
\citep{campello2013hdbscan} clustering. Next, every merge/split decision for deduplication is adjudicated by a single LLM call per candidate group.(2) Result of each stage is checkpointed once it completes, allowing of independent inspection of stages' results and re-using it in follow-up experiments.

\subsection{Initial Typed Triplet Extraction}
\label{sec:extraction}
Subject and object types are extracted as free text, not
constrained to any predefined schema. This schema-free design allows the canonical type vocabulary used throughout the rest of the pipeline to emerge from the data itself.

\subsection{Relation Deduplication}
\label{sec:relation-dedup}

Relation canonicalization maps the raw, free-text \texttt{relation} strings observed across all triplets onto a canonical relation vocabulary
$\mathcal{R}^*$. Candidate clusters are generated by embedding every raw
relation surface form and clustering with HDBSCAN as described
before. Every multi-member candidate cluster is then passed to a single LLM call that decides whether to merge all members into one relation or split them into several. 

\subsection{Entity Types Deduplication}
\label{sec:type-dedup}

This stage both flattens the raw type vocabulary into a canonical set
$\mathcal{T}^*$ and organizes $\mathcal{T}^*$ into a taxonomy hierarchy.

\subsubsection{Type canonicalization}

Candidate clusters over raw type labels are generated the same way as for
relations (\S\ref{sec:relation-dedup}), but with two additional sources of evidence. First, before LLM verification, any two candidate clusters whose relation-argument profiles (TF-IDF-weighted over which canonical relations each type's mentions occur as an argument of) are sufficiently cosine-similar are further merged. Second, each candidate label shown to the LLM verifier is annotated with a the relations it most often co-occurs with, giving the verifier concrete distributional evidence beyond the bare label string.

The verification prompt additionally enforces a hyponymy (``kind of'') test: two labels are merged only if neither is a hyponym of the other, even when their relation-argument profiles overlap almost entirely. For example, \ \emph{documentary film} and \emph{film} participate in the same relations, but LLM instructed to keep them distinct, whereas \emph{movie} and \emph{film}, true synonyms, should 
be merged.

\subsubsection{Hierarchy induction}
\label{sec:hierarchy-induction}

Flat canonicalization removes surface redundancy but does not recover
subsumption structure; a typed knowledge graph needs an explicit is-a
hierarchy over $\mathcal{T}^*$, both to gate candidate pairs during entity
deduplication (\S\ref{sec:entity-dedup}) and to generalize sparse
domain/range observations during constraint induction
(\S\ref{sec:constraints}). We induce this hierarchy with a single
priority-ordered, top-down, pairwise placement procedure:
we evaluated for this step.

% WHY NOT FLAT CLUSTERING?

\paragraph{Priority queue.} Every canonical type in $\mathcal{T}^*$ is
represented by a relation-argument profile, PPMI-weighted by default. A
one-shot pass based on the Weeds precision containment measure
\citep{weeds2004characterising} over the full, untouched pool computes,
for every type $u$, an in-degree score counting how many other types $v$
look containable within it:
\begin{equation}
  \label{eq:weeds-indeg}
  \mathrm{deg_{in}}(u) = \bigl| \{\, v \neq u : \mathrm{WeedsPrec}(v, u) \geq
  \tau \,\} \bigr|.
  \\
\end{equation}
$\mathrm{WeedsPrec}(v, u)$ is the fraction of $v$'s relation-argument weight that $u$
also exhibits:
\begin{equation}
  \mathrm{WeedsPrec}(v, u) = \sum_{f \in F(v)} w_u(f) \big/ \sum_{f \in
F(v)} w_v(f),
\end{equation}
and $F(v)$ is the set of relation-argument features with nonzero
weight under $v$, and $\tau$ is a fixed containment threshold. The in-degree is a structural breadth signal, not a semantic judgment: types are grouped into priority bands from highest in-degree to lowest to identify high-level ontology concept first. If a set of external seed labels is supplied, each is matched to
$\mathcal{T}^*$ by normalized label (or synthesized fresh if unmatched) and pinned at the top of the queue as a root candidates with no other node can be a parent of which; without seeds, the top band of the same ranking performs root discovery instead.

\paragraph{Per-node placement.} Each node is routed top-down through
whatever higher-priority bands have already placed: at every level, the
remaining candidates (initially the whole pool; during descent, an
already-placed node's existing children) are shown to the LLM in
similarity-ranked batches, and a response of ``none of these'' escalates to the next, less-similar batch (up to a fixed cap). Every
batch response is one of four outcomes: 
% \textbf{synonym} of a candidate, \textbf{child of} a candidate, \textbf{parent of} a candidate, \textbf{none} of the candidates.
\begin{itemize}
  \item \textbf{synonym} of a candidate: merge the two synonymous nodes that were not captured in flat clustering deduplication.
  \item \textbf{child of} a candidate: descend one level deeper into that
  candidate's own children at the next iteration, rather than attaching
  immediately, so that a more specific attachment is never preempted.
  \item \textbf{parent of} candidates: insert the incoming node
  \emph{above} the defined candidates, detaching its existing parent edge (if any)
  and reattaching it beneath the new node, which then settles at the
  current level. 
  \item \textbf{none} of the candidates fit at any escalation batch: the
  node settles at its current position -- as a new root if this occurred at
  the top level, or as an ordinary child of its current descent anchor
  otherwise.
\end{itemize}

\paragraph{Depth bound and guarantees.} A maximum-depth cap parameter bounds the worst-case number of placement decisions for any candidate node. A node that would need to exceed it is placed into a structurally separate unresolved bucket rather than being force-attached or discarded, allowing to extent the resulted hieararchy later on. The result of such procedure is always a single-parent forest at every
point in the procedure: acyclicity and the absence of transitive shortcut
edges follow directly from the construction. 

\subsection{Entity Deduplication}
\label{sec:entity-dedup}

Entity deduplication canonicalizes entity mentions and  assigns each canonical entity types. An entity's identity is the compound pair (name, canonical type) rather than name alone, so
that homonymous names typed differently (e.g.\ \emph{Paris} the city versus
\emph{Paris} the person) are never conflated by construction.

Before embedding-based comparison, candidate pairs are restricted by a symbolic gate: two mentions are only ever compared if they share
literally the same type, or are siblings under the same immediate parent
in the induced type hierarchy (\S\ref{sec:hierarchy-induction}). HDBSCAN call only ever compares candidates within one sibling
group: two mentions in different partitions could never have merged under the gate regardless of embedding proximity.

Within each partition, merging proceeds over several rounds rather than a
single pass: each round re-embeds the current pool (already-merged
canonical entities plus untouched singleton mentions), re-clusters with
HDBSCAN, and sends multi-member clusters to the LLM for verification, then folds accepted merges back into the pool for the next round. Rounds stop
once one produces no further reduction in pool size, or a fixed round cap
is reached.

\subsection{Constraints Construction}
\label{sec:constraints}

\paragraph{Stage 0: re-resolution.} For every raw triplet, its subject and
object type strings and its relation string are looked up against the
canonical vocabularies produced by \S\ref{sec:relation-dedup} and
\S\ref{sec:type-dedup}.

\paragraph{Stage 1: relation-direction resolution.} Because relation
canonicalization merges surface forms without regard to grammatical role
(\S\ref{sec:relation-dedup}), a canonical relation can silently pool a
predicate together with its grammatical inverse (e.g.\ \emph{directed} and
\emph{was directed by}), which would otherwise fragment one well-typed
relation's domain/range support across two competing, individually weak
orientations. 

% Subject/object type overlap cannot by itself distinguish this
% case from a genuinely order-free relation with identical domain and range
% types, since both situations produce the same marginal type statistics;
% direction is instead resolved from the signal that actually carries it, the literal surface form. 

For every canonical relation with more than one
surface form, a single LLM call classifies each surface form against the
canonical label as \emph{forward} or \emph{inverted}; triplets whose
original predicate was classified inverted have their subject and object
types swapped before the next stage. This catches merged grammatical
inverses.

\paragraph{Stage 2: hierarchy generalization.} For each canonical
relation's domain (subject-role types) and range (object-role types),
considered independently, we compute the exact lowest common ancestor of
all observed types in the single-parent hierarchy. If the observed types span disconnected
parts of the hierarchy forest, they are partitioned by which root each
descends from and an LCA is computed independently within each partition,
yielding more than one domain (or range) signature and hence more than
one constraint for relations that genuinely connect unrelated type
families.

\paragraph{Stage 3: confidence scoring.} Triplets are regrouped by their
resolved, direction-corrected (domain signature, range signature) pair, and
we compute a partial-completeness-assumption (PCA) style confidence
\citep{galarraga2013amie} for each pair as its joint support over the
relation's total triplet count:
\begin{equation}
  \label{eq:pca-confidence}
  \mathrm{conf}(r, T_d, T_r) = \frac{\mathrm{support}(r, T_d,
  T_r)}{\mathrm{total}(r)}.
\end{equation}
This confidence is thresholded into \textsc{Hard}, \textsc{Soft}, and
\textsc{Hint} constraint strengths.

% \section{KG Retrieval and QA}
% TODO -- not yet written.

% \section{Evaluation}
% TODO -- not yet written.

\section*{Limitations}
% TODO

\section*{Acknowledgments}
% TODO

\bibliography{custom}

\end{document}
